\section{Task Definition}
\Secl{task}

Each problem in the \FVSpec benchmark consists of four components: (1) the original Python implementation of the function under test, (2) the original Python PBT exercising that implementation, (3) a full Lean implementation of the same function, and (4) a Lean theorem corresponding to the PBT, with a \texttt{sorry} placeholder where the proof should go. The AI's task is to fill in every theorem \texttt{sorry} with a Lean proof term that compiles cleanly under \texttt{lake build}; submissions are then scored by whether the proof retains any \texttt{sorry} or special-purpose axioms. Each problem also includes a URL pointing to the original source code and license information.

As a running example, consider \texttt{isort(lst)}. Figure~\ref{fig:running-example}
shows all four artifacts of a single \FVSpec problem. The left column holds
the original implementation (top) and the PBT exercising it
(bottom); the right column holds the corresponding \texttt{Impl.lean}
(top), a complete Lean port of the function, and \texttt{Spec.lean}
(bottom), containing two theorems---one for sortedness, one for permutation---
each with their own \texttt{sorry}. The AI's task is to discharge both
Lean \texttt{sorry} placeholders.

\begin{lrbox}{\fvspecPyImplBox}%
\begin{minipage}{0.36\linewidth}
{\small\sffamily\textbf{Python: implementation of \texttt{isort}}}\par\smallskip
\begin{lstlisting}[style=code, language=Python, basicstyle=\ttfamily\scriptsize]
from bisect import insort

def isort(lst):
    out = []
    for x in lst: insort(out, x)
    return out
\end{lstlisting}
\end{minipage}%
\end{lrbox}%

\begin{lrbox}{\fvspecImplBox}%
\begin{minipage}{0.46\linewidth}
{\small\sffamily\textbf{Lean: \texttt{Impl.lean}}}\par\smallskip
\begin{lstlisting}[style=code, language=Lean, basicstyle=\ttfamily\scriptsize]
namespace Fvspec.Impl

def insert (x : Int) : List Int -> List Int
  | []      => [x]
  | y :: ys => if x <= y then x :: y :: ys
               else y :: insert x ys

def isort : List Int -> List Int
  | []      => []
  | x :: xs => insert x (isort xs)

end Fvspec.Impl
\end{lstlisting}
\end{minipage}%
\end{lrbox}%

\begin{lrbox}{\fvspecPyPbtBox}%
\begin{minipage}{0.36\linewidth}
{\small\sffamily\textbf{Python: Hypothesis PBT}}\par\smallskip
\begin{lstlisting}[style=code, language=Python, basicstyle=\ttfamily\scriptsize]
from collections import Counter
from hypothesis import given
import hypothesis.strategies as st
@given(st.lists(st.integers()))
def test_isort(lst):
    out = isort(lst)
    pairs = zip(out, out[1:])
    assert all(a <= b for a, b in pairs)
    assert Counter(out) == Counter(lst)
\end{lstlisting}
\end{minipage}%
\end{lrbox}%

\begin{lrbox}{\fvspecSpecBox}%
\begin{minipage}{0.46\linewidth}
{\small\sffamily\textbf{Lean: \texttt{Spec.lean}}}\par\smallskip
\begin{lstlisting}[style=code, language=Lean, basicstyle=\ttfamily\scriptsize]
import Fvspec.Impl
open Fvspec.Impl

theorem isort_sorted (l : List Int) :
  (isort l).Sorted LE.le := sorry

theorem isort_perm (l : List Int) :
  (isort l).Perm l := sorry
\end{lstlisting}
\end{minipage}%
\end{lrbox}%

\begin{figure}[h]
\centering
\begin{tikzpicture}[
  box/.style={draw, rounded corners, inner sep=4pt, fill=gray!5},
  arr/.style={-{Latex[length=2.5mm]}, thick, gray!50!black},
  group/.style={draw, dashed, rounded corners, inner sep=2mm,
                gray!60!black, line width=0.6pt}
]
  \node[box, anchor=east] (pyimpl)
       at (-0.5cm,  2.2cm) {\usebox{\fvspecPyImplBox}};
  \node[box, anchor=west] (limpl)
       at ( 0.5cm,  2.2cm) {\usebox{\fvspecImplBox}};
  \node[box, anchor=east] (pypbt)
       at (-0.5cm, -2.2cm) {\usebox{\fvspecPyPbtBox}};
  \node[box, anchor=west] (lspec)
       at ( 0.5cm, -2.2cm) {\usebox{\fvspecSpecBox}};
  \node[group, fit=(pyimpl)(pypbt),
        label={[font=\small\sffamily\bfseries]above:{\FVSpec:PBT}}] {};
  \node[group, fit=(limpl)(lspec),
        label={[font=\small\sffamily\bfseries]above:{\FVSpec:FV}}] {};
  \draw[arr] (pyimpl.east) -- (limpl.west);
  \draw[arr] (pypbt.east)  -- (lspec.west);
\end{tikzpicture}
\caption{The four artifacts of a single \FVSpec problem. Each problem is
partitioned into the \emph{PBT corpus} (\FVSpec:PBT, left)---the original
Python implementation and Hypothesis test we scrape from GitHub---and the
\emph{formal-verification challenge} (\FVSpec:FV, right) we transpile from
it: a complete Lean port of the implementation and a Lean theorem statement
of the property under test. The AI must discharge both Lean \texttt{sorry}
placeholders, one per theorem.}
\label{fig:running-example}
\end{figure}

% Next, we explain how we scrape the PBT corpus and what the corpus looks like.
